% ======================================================================
% Ontologie der Kontinua -- Core 1.3
% Cross-K-Modul: crossk_k4_k5.tex
% Cross-Level-Relationen zwischen K4 und K5
% ======================================================================

\subsubsection{K4–K5-Querverbindung}
\label{sec:crossk-k4-k5}

Der Übergang $K_4 \rightarrow K_5$ markiert die Entstehung des neuronalen
Kontinuums aus dem biologischen Kontinuum. Während $K_4$ Gradienten,
Membranen und metabolische Zyklen unterstützt, führt $K_5$ ein:
\begin{itemize}
  \item elektrische Erregbarkeit und Membranspannungsdynamik,
  \item Ionenkanäle mit diskreten Öffnungs-/Schließzuständen,
  \item propagierende elektrische Ereignisse (Proto-Spikes),
  \item Refraktärmechanismen und Erregbarkeitsschwellen,
  \item räumliche Ausbreitung entlang von Membranen und frühe Netzwerkstruktur,
  \item attractorartige Muster, die Proto-Informationsflüsse bilden.
\end{itemize}

Dies ist die Geburt einer nervensystemartigen Organisation: Struktur, die elektrische
Signale kodieren, propagieren und transformieren kann.

% ----------------------------------------------------------------------
\subsubsection{1. Stufen und ihre Rollen}

\paragraph{K4 (biologisches Kontinuum).}
$K_4$ liefert:
\begin{itemize}
  \item Membranen und Kompartimente $\partial\Omega$,
  \item Gradienten (Ionen, pH, Redox) mit Potenzialen $P_{\mathrm{grad}}$,
  \item Ionenfluss und Osmose ($J_{\mathrm{ion}}, J_{\mathrm{osm}}$),
  \item metabolische und Pumpzyklen $C_{\mathrm{energy}}$,
  \item Schwellen $\Theta_{\mathrm{mem}}, \Theta_{\mathrm{grad}},
        \Theta_{\mathrm{osm}}, \Theta_{\mathrm{redox}}$.
\end{itemize}

Allerdings fehlen in $K_4$:
\begin{itemize}
  \item regenerative elektrische Erregung,
  \item diskrete Kanalgating-Dynamik,
  \item großräumige elektrische Ausbreitung,
  \item stabile elektrische Attraktoren.
\end{itemize}

\paragraph{K5 (neuronales Kontinuum).}
$K_5$ führt ein:
\begin{itemize}
  \item die Erregungsachse $A_{\mathrm{exc}}$ als Membranspannung,
  \item Achsen der Ionenkanäle $A_{\mathrm{channel}}$ mit Öffnungs-/Schließzuständen,
  \item elektrische Potentiale $P_{\mathrm{elect}}$,
  \item Flüsse $J_5 = (J_{\mathrm{ion}}, J_{\mathrm{exc}}, J_{\mathrm{leak}},
        J_{\mathrm{shunt}})$,
  \item Proto-Spike- und Refraktärzyklen $C_{\mathrm{spike}}, C_{\mathrm{recovery}}$,
  \item Schwellen $\Theta_{\mathrm{exc}}, \Theta_{\mathrm{refrac}},
        \Theta_{\mathrm{noise-electrical}}$.
\end{itemize}

$K_5$ ist das erste Kontinuum, das Spatio-zeitliche Musterausbreitung
leisten kann.

% ----------------------------------------------------------------------
\subsubsection{2. Vererbte und neue Achsen sowie Schwellen}

\paragraph{Ererbte Achsen.}
Die Gradienten von $K_4$ werden zu Komponenten der elektrischen Achse:
\[
A_{\mathrm{grad}}(K_4) \subset A_{\mathrm{exc}}(K_5).
\]
Die Membranachse $A_{\mathrm{mem}}$ bildet den Untergrund für Kanalverteilung und
räumliche Ausbreitung.

\paragraph{Neue Achsen.}
$K_5$ führt ein:
\begin{itemize}
  \item $A_{\mathrm{exc}}$ — Achse der Membranspannung,
  \item $A_{\mathrm{channel}}$ — diskrete Gating-Zustände,
  \item $A_{\mathrm{lat}}$ — laterale Ausbreitung entlang der Membranoberfläche.
\end{itemize}

\paragraph{Schwellen.}
Neue Schwellenfamilien umfassen:
\begin{itemize}
  \item $\Theta_{\mathrm{exc}}$ — minimale Depolarisation zur Spike-Entstehung,
  \item $\Theta_{\mathrm{refrac}}$ — Erholungszeitbeschränkungen,
  \item $\Theta_{\mathrm{channel}}$ — Stabilität von Gating und Versagenszuständen,
  \item $\Theta_{\mathrm{noise-electrical}}$ — rauschinduzierte Kollapsgrenzen.
\end{itemize}

Verletzung von $\Theta_4$ verhindert Erregbarkeit:
\[
\Theta_4^{\mathrm{death}} \Rightarrow \Omega(K_5)=\varnothing.
\]

% ----------------------------------------------------------------------
\subsubsection{3. Cross-Level-Flüsse, Zyklen und Spannungen}

\paragraph{Flüsse.}
Flüsse in $K_5$ leiten sich aus $J_4$ ab, enthalten aber neue elektrische
Komponenten:
\[
J_5 = (J_{\mathrm{ion}}, J_{\mathrm{exc}}, J_{\mathrm{leak}}, J_{\mathrm{shunt}}).
\]

Wichtige Beiträge:
\begin{itemize}
  \item \textbf{$J_{\mathrm{exc}}$}: regenerative Natrium-ähnliche Entladung,
  \item \textbf{$J_{\mathrm{leak}}$}: dämpfender Hintergrundstrom,
  \item \textbf{$J_{\mathrm{shunt}}$}: membranstabilisierende Bypass-Ströme,
  \item \textbf{$J_{\mathrm{ion}}$}: klassischer Nernst-ähnlicher Gradientenfluss.
\end{itemize}

\paragraph{Zyklen.}
Zwei fundamentale Zyklen entstehen in $K_5$:
\[
C_{\mathrm{spike}}:\ \Delta V \uparrow \rightarrow \Delta V_{\mathrm{peak}}
\rightarrow \text{Repolarisation} \rightarrow \text{Erholung}.
\]

\[
C_{\mathrm{recovery}}:\ \text{Refraktär} \rightarrow \text{Reset}
\rightarrow \text{Erregbarkeit hergestellt}.
\]

Diese Zyklen hängen von stabilen Gradienten und Kanalkinetik ab.

\paragraph{Spannung.}
Cross-Level-Spannung:
\[
T_{4,5} =
g(\Theta_{\mathrm{exc}} - P_{\mathrm{elect}},\
  J_{\mathrm{leak}} - J_{\mathrm{exc}},\
  \Theta_{\mathrm{channel}} - A_{\mathrm{channel}}).
\]

Wenn $T_{4,5}$ die dimensionsbezogene Schwelle überschreitet, bricht die
Propagation ab und das System kollabiert in ein nicht-erregbares K₄-Regime.

% ----------------------------------------------------------------------
\subsubsection{4. Geburts- und Todesbedingungen über die Stufen}

\paragraph{Geburt von \texorpdfstring{$K_5$}{K_5}.}
Das neuronale Kontinuum entsteht, wenn:
\[
T_4 > \Theta_{4,\mathrm{dim}}
\quad\text{und}\quad
A_{\mathrm{exc}}, A_{\mathrm{channel}} \in M_5 \setminus A(K_4),
\]
was stabile Spike-Dynamik erlaubt.

Erweiterung des Zustandsraums:
\[
\Omega(K_5) = \Omega(K_4) \cup \Delta\Omega_{\mathrm{neural}}.
\]

\paragraph{Todesausbreitung.}
Versagt $\Theta_4$ (Membranruptur, Gradientenkollaps), verliert $K_5$ die Basis
für Erregbarkeit.

Versagt $\Theta_5$ (zu viel Rauschen, Instabilität der Kanäle,
unzureichender Gradient), kollabiert nur $K_5$; $K_4$ bleibt erhalten.

% ----------------------------------------------------------------------
\subsubsection{5. Konzeptionelle Beispiele}

\paragraph{Proto-Spike-Bildung.}
Gegeben:
\[
\Delta V > \Theta_{\mathrm{exc}}
\quad\text{und}\quad
A_{\mathrm{channel}} = \text{offen},
\]
entsteht eine regenerative Aufstrampelfunktion und propagiert über die lokale
Membran.

\paragraph{2D-Propagation.}
Die seitliche Achse $A_{\mathrm{lat}}$ ermöglicht Ausbreitung elektrischer Muster
über zweidimensionale Membranoberflächen — unmöglich in $K₄$.

\paragraph{Erregbarkeit als neue Dimension.}
Das Auftreten von $\Delta V(t)$-Dynamik bildet eine neue Achse und eine neue
Klasse von Flüssen, die die dimensionsbezogene Wachstumsregel erfüllen:
\[
\dim K_5 = \dim K_4 + 1.
\]

\paragraph{Frühe Informationsflüsse.}
Spike-ähnliche Aktivität unterstützt Musterübertragung und Proto-Kodierung und
erzeugt die erste Grundlage für strukturierte Information innerhalb der
Kontinuums-Hierarchie.

Somit lässt sich der K4→K5-Übergang in komprimierter Form zusammenfassen als die
Entstehung von Erregbarkeit, elektrischer Signalisierung und Musterpropagation aus
biologischen Gradienten und Membranen.
